\chapter{Réalisation et résultats}

\section*{Introduction}
\addcontentsline{toc}{section}{Introduction}

Ce chapitre rend compte de la réalisation effective du projet sur le VPS de la BTE.
Il s'ouvre sur la description de l'environnement matériel et logiciel sur lequel la
plateforme a été déployée, puis enchaîne sur l'implémentation sprint par sprint
selon la méthodologie Scrumban introduite au chapitre 1. Chaque sprint suit le même
gabarit (objectif inscrit au backlog, obstacles rencontrés en production, livrable
effectif). Ce choix de présentation laisse apparaître l'écart entre le plan prévu
et ce que la production a effectivement renvoyé. La dernière partie présente les
résultats obtenus sur des \textit{Pull Requests} réelles ainsi que la validation de
bout en bout de la plateforme.

% ─────────────────────────────────────────────────────────────────────────────
\section{Environnement de travail}

    \subsection{Environnement matériel}

        Le tableau \ref{tab:vps_env} présente les caractéristiques du VPS de production
        sur lequel l'ensemble de la plateforme a été déployé.

        \begin{table}[H]
        \centering
        \begin{tabular}{|m{3cm}|m{3cm}|m{6cm}|}
        \hline
        \textbf{Composant} & \textbf{Valeur} & \textbf{Détail} \\
        \hline
        CPU & 12 cœurs & Intel Haswell, AVX2 \\
        \hline
        RAM & 45 Go & \textasciitilde41 Go disponibles \\
        \hline
        Disque & 290 Go & \textasciitilde245 Go libres \\
        \hline
        GPU & Aucun & CPU-only AVX2 \\
        \hline
        OS & Ubuntu Linux & kernel 6.14.0-37 \\
        \hline
        \end{tabular}
        \caption{Environnement matériel — VPS de production BTE}
        \label{tab:vps_env}
        \end{table}

    \subsection{Environnement logiciel}

        Les versions des composants logiciels déployés sont celles arrêtées à l'issue
        de l'étude technologique du chapitre~2 et récapitulées dans le tableau
        \ref{tab:tech_synthesis} : Docker Engine 29.4.0 \cite{Docker}, Python 3.12,
        FastAPI \cite{FastAPI}, LangGraph \cite{LangGraph}, PostgreSQL 16
        \cite{PostgreSQL}, Redis 7 \cite{Redis} et les cinq scanners SAST
        \cite{Trivy, Gitleaks, Semgrep, Checkov, OSVScanner}.

% ─────────────────────────────────────────────────────────────────────────────
\section{Implémentation par sprints Scrumban}

    La réalisation s'est étalée sur dix itérations : huit sprints majeurs
    (Sprint~1 à Sprint~8) qui ont construit la plateforme dans sa forme livrée,
    et deux micro-sprints d'extension (Post-Sprint~8 et Sprint~8.1) menés après la
    clôture du périmètre initial pour évaluer un second backend d'inférence et
    étendre la couverture d'observabilité.

    % Chaque sprint est présenté ci-dessous
    %selon le même encart standardisé (\textit{Objectif du sprint}, \textit{Obstacles
    %rencontrés}, \textit{Livrable}), suivi du détail technique des décisions, des
    %configurations et des correctifs déployés.

    \subsection{Sprint 1 — Fondations : VPS, Docker et webhook GitHub}

        \begin{description}
            \itemsep 0.3em
            \item[Objectif du sprint :] Mettre en place l'ensemble de l'infrastructure
                Docker, configurer Ollama pour un fonctionnement CPU-only et valider
                la réception du premier webhook GitHub.
            \item[Obstacles rencontrés :] Nginx, configuré en HTTP/1.0 par défaut,
                altérait les corps \textit{chunked} des requêtes GitHub, provoquant
                systématiquement des erreurs 500 en réception.
            \item[Livrable :] Onze conteneurs déployés et opérationnels, premier
                webhook GitHub reçu et validé par signature HMAC-SHA256, configuration
                Ollama tunée pour CPU Haswell AVX2.
        \end{description}

        Le déploiement initial (installation de Docker Engine sur l'Ubuntu existant,
        règles \texttt{iptables} limitant l'exposition externe aux ports 80 et 443,
        initialisation du schéma PostgreSQL, démarrage des onze conteneurs) relève de
        l'administration système courante et n'est pas détaillé ici. Deux réglages
        méritent en revanche d'être exposés.

        Le premier est la configuration de performance d'Ollama, déterminante en
        fonctionnement CPU-only : avec les paramètres par défaut, le modèle 14B
        consomme plus de onze Go de mémoire KV à seize mille tokens de contexte. La
        figure \ref{fig:ollama_compose} présente la configuration retenue dans le
        fichier \texttt{docker-compose.\allowbreak{}yml}.

        \begin{figure}[H]
            \centering
            \frame{\includegraphics[width=0.88\columnwidth]{ollama_docker_compose}}
            \caption{Configuration Ollama dans docker-compose.yml}
            \label{fig:ollama_compose}
        \end{figure}

        La variable \texttt{OLLAMA\_\allowbreak{}FLASH\_\allowbreak{}ATTENTION=1} réduit la complexité mémoire du
        cache KV de O(n²) à O(n) et économise environ 650 Mo à seize mille tokens de
        contexte. La variable \texttt{OLLAMA\_\allowbreak{}KV\_\allowbreak{}CACHE\_\allowbreak{}TYPE=q8\_\allowbreak{}0} applique une
        quantisation huit bits au cache KV, divisant sa consommation mémoire par deux.
        La mémoire partagée est fixée à deux Go pour les buffers de synchronisation
        des threads, contre une valeur par défaut de soixante-quatre Mo insuffisante
        pour douze cœurs parallèles. La limite mémoire de quarante-deux Go protège les
        autres services du VPS tout en permettant au modèle 14B (neuf Go) de fonctionner
        avec une marge suffisante.

        Le second est le correctif du problème de transport signalé en tête de sprint.
        La figure \ref{fig:nginx_fix} présente la configuration nginx corrigée.

        \begin{figure}[H]
            \centering
            \frame{\includegraphics[width=0.88\columnwidth]{nginx_webhook_fix}}
            \caption{Correctif nginx}
            \label{fig:nginx_fix}
        \end{figure}

        L'ajout de \texttt{proxy\_\allowbreak{}http\_\allowbreak{}version 1.\allowbreak{}1} et \texttt{proxy\_\allowbreak{}set\_\allowbreak{}header Connection ""}
        a résolu ce problème en forçant nginx à utiliser des connexions persistantes
        HTTP/1.1 avec l'agent, ce qui met fin à l'altération des corps \textit{chunked}.

        \paragraph{Configuration de l'intégration GitHub:}
        Côté GitHub, deux configurations ont été nécessaires pour permettre la
        communication bidirectionnelle entre la plateforme et le dépôt cible.
        En premier lieu, un \textit{Personal Access Token} a été créé sur le
        compte de l'agent avec les scopes \texttt{repo} (lecture/écriture du
        contenu) et \texttt{write:discussion} (publication des commentaires
        \textit{inline} et mise à jour du statut du commit), comme présenté
        sur la figure \ref{fig:github_token_creation}. Ce token est injecté
        dans l'agent via la variable d'environnement \texttt{GITHUB\_\allowbreak{}TOKEN}
        et stocké en \textit{secret} Docker pour ne pas apparaître dans les
        logs ou les images.

        \begin{figure}[H]
            \centering
            \frame{\includegraphics[width=0.88\columnwidth]{github_token_creation}}
            \caption{Création du Personal Access Token GitHub}
            \label{fig:github_token_creation}
        \end{figure}

        En second lieu, un \textit{webhook} a été configuré sur le dépôt cible
        pour notifier l'agent à chaque événement de \textit{Pull Request}.
        La figure \ref{fig:github_webhook_config} présente cette configuration :
        l'URL pointe vers l'endpoint public \texttt{/\allowbreak{}webhooks/\allowbreak{}github} de l'agent,
        un secret partagé est défini pour la signature HMAC-SHA256 des charges
        utiles, et le \textit{content type} est fixé à \texttt{application/\allowbreak{}json}
        conformément à ce qu'attend le router FastAPI.

        \begin{figure}[H]
            \centering
            \frame{\includegraphics[width=0.88\columnwidth]{github_webhook_config}}
            \caption{Configuration du \textit{webhook} GitHub sur le dépôt cible}
            \label{fig:github_webhook_config}
        \end{figure}

        La figure \ref{fig:webhook_received_github} confirme la réception correcte
        du webhook côté GitHub après ce correctif : l'événement \textit{pull
        request} est livré avec succès à l'endpoint de l'agent, qui répond par un
        code 202 \textit{Accepted} validant la prise en charge asynchrone de la
        requête.

        \begin{figure}[H]
            \centering
            \frame{\includegraphics[width=0.88\columnwidth]{webhook_received_github}}
            \caption{Webhook GitHub correctement reçu par l'agent}
            \label{fig:webhook_received_github}
        \end{figure}

        \paragraph{Traitement applicatif de l'événement:}
        La livraison côté GitHub ne suffit pas : il faut aussi décrire ce que
        l'agent fait une fois la requête reçue. La figure
        \ref{fig:github_webhook_handler} présente le handler implémenté dans
        \texttt{app/\allowbreak{}routers/\allowbreak{}webhooks.\allowbreak{}py}. L'endpoint
        \texttt{POST /\allowbreak{}webhooks/\allowbreak{}github} répond en 202 et confie le reste à une
        tâche de fond, afin de rester dans les délais de réponse imposés par
        GitHub.

        \begin{figure}[H]
            \centering
            \frame{\includegraphics[width=0.92\columnwidth]{github_webhook_handler_code}}
            \caption{Handler FastAPI de réception des webhooks GitHub}
            \label{fig:github_webhook_handler}
        \end{figure}

        Le traitement commence par la vérification de la signature
        HMAC-SHA256 : le corps brut de la requête est comparé à l'en-tête
        \texttt{X-Hub-Signature-256}. Une signature absente ou incorrecte
        entraîne un rejet immédiat en 403, avant tout décodage JSON. Le corps
        est ensuite interprété selon le \textit{content type} reçu. GitHub
        peut transmettre la charge utile en JSON direct, ou l'encapsuler dans
        un champ \texttt{payload} au format
        \texttt{application/\allowbreak{}x-www-form-urlencoded}. Cette seconde branche a
        été ajoutée après les premiers tests d'intégration, lorsque certaines
        livraisons ne correspondaient pas au format JSON initialement prévu.

        Seuls les événements \texttt{pull\_\allowbreak{}request} dont l'action vaut
        \texttt{opened} ou \texttt{synchronize} sont conservés. Les autres
        retournent une réponse HTTP 200 avec un message d'ignorance, sans
        lancer de pipeline. Pour un événement accepté, un identifiant de
        tâche est généré et la fonction \texttt{dispatch\_\allowbreak{}event} est
        enregistrée via \texttt{BackgroundTasks} de FastAPI. L'appel HTTP se
        termine sans attendre la fin de l'analyse ; la déduplication Redis et
        le lancement du graphe LangGraph sont pris en charge par le
        dispatcher, déjà présenté au chapitre~3.

    \subsection{Sprint 2 — Pipeline de sécurité : scanners et LLM}

        \begin{description}
            \itemsep 0.3em
            \item[Objectif du sprint :] Intégrer les cinq scanners de sécurité
                (Trivy, Gitleaks, Semgrep, Checkov, OSV-Scanner) et le modèle LLM 14B
                dans le workflow LangGraph.
            \item[Obstacles rencontrés :] Le diff retourné par l'API GitHub avec ses
                trois lignes de contexte par défaut s'est révélé trop pauvre pour
                permettre au modèle d'identifier le contexte d'une vulnérabilité.
            \item[Livrable :] Génération d'un diff local avec quinze lignes de
                contexte, exécution parallèle des cinq scanners via
                \texttt{asyncio.\allowbreak{}gather()} avec isolation des erreurs par coroutine,
                mise en cache Redis des résultats avec TTL d'une heure.
        \end{description}

        Le diff transmis au modèle est désormais généré localement par
        \texttt{git diff -U15}, avec quinze lignes de contexte au lieu des trois que
        l'API GitHub renvoie par défaut. La figure \ref{fig:diff_comparison} illustre
        la différence entre les deux approches sur un exemple d'injection SQL : avec
        trois lignes seulement, le modèle ne voit ni l'origine de la variable
        concaténée ni la fonction qui l'entoure.

        \begin{figure}[H]
            \centering
            \frame{\includegraphics[width=0.88\columnwidth]{diff_comparison}}
            \caption{Comparaison diff API GitHub (-U3) vs diff local (-U15)}
            \label{fig:diff_comparison}
        \end{figure}

        L'exécution parallèle des scanners via \texttt{asyncio.\allowbreak{}gather()} est un élément
        clé des performances du pipeline. Chaque scanner s'exécute dans sa propre coroutine
        avec isolation des erreurs, et son résultat est mis en cache dans Redis. La figure
        \ref{fig:scan_parallel} montre le résultat de cette exécution parallèle avec les
        sorties des cinq scanners, ici sur la PR~\#18 du dépôt
        \texttt{GhaiethFerchichi/\allowbreak{}Vunl-application}.

        \begin{figure}[H]
            \centering
            \frame{\includegraphics[width=0.88\columnwidth]{scanners_parallel_output}}
            \caption{PR~\#18 — exécution parallèle des cinq scanners}
            \label{fig:scan_parallel}
        \end{figure}

    \subsection{Sprint 3 — Revue combinée et commentaires inline GitHub}

        \begin{description}
            \itemsep 0.3em
            \item[Objectif du sprint :] Réduire la durée du pipeline et publier les
                commentaires inline directement sur les fichiers modifiés de la
                \textit{Pull Request}.
            \item[Obstacles rencontrés :] L'analyse de sécurité et la revue de
                qualité de code nécessitaient deux appels successifs au modèle 14B
                totalisant entre trente et cinquante minutes, durée incompatible
                avec la fluidité d'un pipeline CI/CD.
            \item[Livrable :] Fusion des deux appels 14B en un seul appel combiné
                dans le nœud \texttt{analyze\_\allowbreak{}review\_\allowbreak{}node}, durée totale du pipeline
                ramenée à quinze à vingt-cinq minutes (gain d'environ 50\,\% par rapport
                aux deux appels séquentiels), commentaires inline opérationnels avec
                bouton \textit{Apply}.
        \end{description}

        Le gain annoncé porte sur la durée totale du pipeline (scanners,
        classification 7B et publication GitHub comprises), et non sur la seule phase
        d'inférence du modèle 14B. Le tableau \ref{tab:combined_perf} résume cet
        impact.

        \begin{table}[H]
        \centering
        \begin{tabular}{|m{5cm}|m{3cm}|m{3cm}|m{2cm}|}
        \hline
        \textbf{Configuration} & \textbf{Avant fusion (Sprint 3)} & \textbf{Après fusion (Sprint 3)} & \textbf{Gain} \\
        \hline
        Appels LLM 14B par PR & 2 séquentiels & 1 combiné & -1 appel \\
        \hline
        Durée totale du pipeline & 30–50 min & 15–25 min & $\sim$50\,\% \\
        \hline
        \texttt{num\_ctx} utilisé & 2 × 8 192 & 12 288 & +contexte \\
        \hline
        \end{tabular}
        \caption{Impact de la fusion des appels LLM sur la performance du pipeline}
        \label{tab:combined_perf}
        \end{table}

        La PR~\#18 sert de fil rouge pour valider ce sprint. Le fichier
        \texttt{banking\_\allowbreak{}demo\_\allowbreak{}final.\allowbreak{}php} y regroupe cinq failles types :
        injection SQL, secrets en dur, path traversal, injection de commande
        et hachage faible. La revue de sécurité publiée sur GitHub est reproduite
        intégralement dans la section consacrée aux résultats
        (figures~\ref{fig:pr_full_result} à~\ref{fig:pr_full_result_3}) ; on se
        limite ici aux trois sorties propres à ce sprint. La
        figure~\ref{fig:github_review_quality} montre la revue de qualité de code
        publiée dans l'onglet \textit{Conversation}.

        \begin{figure}[H]
            \centering
            \frame{\includegraphics[width=0.92\columnwidth]{pr_comment-code_quality_review_0}}
            \caption{PR~\#18 — revue de qualité de code publiée dans l'onglet \textit{Conversation}.}
            \label{fig:github_review_quality}
        \end{figure}

        La figure~\ref{fig:inline_suggestions} reprend la même PR dans
        l'onglet \textit{Files changed} : cinq commentaires \textit{inline}
        avec suggestion de correction et bouton \textit{Apply}.

        \begin{figure}[H]
            \centering
            \frame{\includegraphics[width=0.92\columnwidth]{github_inline_suggestions}}
            \caption{PR~\#18 — suggestion \textit{inline} sur l'injection SQL}
            \label{fig:inline_suggestions}
        \end{figure}

        Côté serveur, la figure~\ref{fig:agent_logs_pr_complete} montre le
        nœud \texttt{report\_node} au travail : revue postée, cinq suggestions
        \textit{inline}, statut commit passé à \texttt{error}, notification
        Slack envoyée. Durée totale : environ quinze minutes pour cette PR.

        \begin{figure}[H]
            \centering
            \frame{\includegraphics[width=0.92\columnwidth]{agent_logs_pr_complete}}
            \caption{PR~\#18 — journaux agent lors de la finalisation}
            \label{fig:agent_logs_pr_complete}
        \end{figure}

        La figure~\ref{fig:slack_notification_pr_result} reprend le même
        résultat dans le canal \texttt{\#bte-security-ai-agent} :
        dépôt, numéro de PR, risque \texttt{HIGH}, verdict \texttt{BLOCK}.

        \begin{figure}[H]
            \centering
            \frame{\includegraphics[width=0.80\columnwidth]{slack_notification_pr_result}}
            \caption{PR~\#18 — notification Slack de fin de revue}
            \label{fig:slack_notification_pr_result}
        \end{figure}

    \subsection{Sprint 4 — Opérations autonomes et urgence disque}

        \begin{description}
            \itemsep 0.3em
            \item[Objectif du sprint :] Rendre la plateforme auto-réparante via une
                surveillance disque autonome et un \textit{health digest} quotidien
                publié sur Slack.
            \item[Obstacles rencontrés :] Incident critique du 20 avril 2026. Un
                téléchargement partiel du modèle \texttt{qwen2.\allowbreak{}5-coder:32b} a laissé
                un blob de 242~Go sur le disque sans manifeste associé, portant
                l'utilisation à 92\,\%.
            \item[Livrable :] 233~Go récupérés sans aucune interruption de service,
                gardien disque autonome surveillant l'utilisation toutes les trente
                minutes, \textit{scheduler} asyncio avec \textit{health digest}
                quotidien publié à 09h00 UTC.
        \end{description}

        \subsubsection{Urgence disque du 20 avril 2026}

        L'incident date du sprint trois ; nous le racontons ici parce que c'est
        lui qui a poussé à coder le gardien disque de ce sprint. Un téléchargement
        interrompu du modèle \texttt{qwen2.\allowbreak{}5-coder:32b} avait laissé un blob de
        242~Go sans manifeste. Avec le cache Docker et le modèle
        \texttt{llama3.2:3b} toujours présent, le disque est monté à 92\,\%.
        Nous n'avons pas gardé de capture d'écran de la session de dépannage ;
        les journaux du VPS et une carte Jira en classe \textit{Expedite} suffisent
        à la retrouver. Environ 233~Go ont été libérés sans arrêter les conteneurs.

        Le gardien disque décrit à la figure~\ref{fig:scheduler} est la réponse
        directe à cet épisode. L'état stable du disque après reprise apparaît
        dans le digest Slack de la figure~\ref{fig:slack_digest}.

        \subsubsection{Scheduler autonome et health digest}

        La figure \ref{fig:scheduler} présente le scheduler autonome avec ses deux tâches
        asyncio permanentes : le gardien disque et le health digest quotidien.

        \begin{figure}[H]
            \centering
            \frame{\includegraphics[width=0.65\columnwidth,height=0.80\textheight,keepaspectratio]{scheduler_tasks}}
            \caption{Scheduler autonome}
            \label{fig:scheduler}
        \end{figure}

        Pour permettre à l'agent d'émettre des notifications sur Slack, un
        \textit{bot} a été créé dans l'espace \texttt{team-devsecops} avec les
        portées \texttt{channels:read}, \texttt{chat:write} et
        \texttt{chat:write.\allowbreak{}public}, comme le montre la
        figure~\ref{fig:slack_bot_config}. Son \textit{bot token} est injecté
        via \texttt{SLACK\_\allowbreak{}BOT\_\allowbreak{}TOKEN}, l'identifiant du canal via
        \texttt{SLACK\_\allowbreak{}CHANNEL\_\allowbreak{}ID}.

        \begin{figure}[H]
            \centering
            \frame{\includegraphics[width=0.88\columnwidth]{slack_bot_config}}
            \caption{Configuration OAuth du \textit{bot} Slack}
            \label{fig:slack_bot_config}
        \end{figure}

        La figure \ref{fig:slack_digest} montre un exemple du health digest Slack publié
        quotidiennement à 09h00 UTC, récapitulant l'état du disque, de tous les conteneurs,
        du modèle Ollama chargé et des alertes Prometheus actives.

        \begin{figure}[H]
            \centering
            \frame{\includegraphics[width=0.70\columnwidth]{slack_health_digest}}
            \caption{Health digest quotidien Slack}
            \label{fig:slack_digest}
        \end{figure}

    \subsection{Sprint 5 — Stack de monitoring Prometheus et Grafana}

        \begin{description}
            \itemsep 0.3em
            \item[Objectif du sprint :] Déployer la stack Prometheus + Grafana avec
                trois dashboards opérationnels pour la surveillance de la plateforme.
            \item[Obstacles rencontrés :] Au premier démarrage, les trois cibles
                Prometheus affichaient toutes l'état \texttt{down}. L'option
                \texttt{--web.\allowbreak{}route-prefix} de Prometheus et AlertManager préfixait
                l'ensemble de leurs chemins HTTP, ce qui rendait les chemins de
                \textit{scrape} par défaut incorrects.
            \item[Livrable :] Quatre cibles Prometheus en état \texttt{up} après
                ajout des \texttt{metrics\_\allowbreak{}path} explicites, trois dashboards Grafana
                provisionnés (VPS Host Monitoring, DevSecOps AI Agent, PR Security
                Reviews).
        \end{description}

        \subsubsection{Correction des métriques Prometheus}

        Lors du premier démarrage de la stack Prometheus \cite{Prometheus},
        les trois cibles configurées affichaient l'état \texttt{down}. Prometheus
        et AlertManager \cite{AlertManager} tournent avec
        \texttt{--web.\allowbreak{}route-prefix}, ce qui déplace le chemin \texttt{/metrics}
        par rapport à ce que le scrape attend par défaut. Nous n'avons pas
        archivé la page \textit{Targets} dans cet état initial ; le diagnostic
        repose sur les journaux de scrape et la lecture de
        \texttt{prometheus.\allowbreak{}yml}.

        Après ajout des \texttt{metrics\_\allowbreak{}path} explicites (\texttt{/\allowbreak{}prometheus/\allowbreak{}metrics}
        pour Prometheus et \texttt{/\allowbreak{}alertmanager/\allowbreak{}metrics} pour AlertManager), ainsi que
        l'ajout de la cible node-exporter, les quatre cibles sont passées à l'état
        \texttt{up}, comme le montre la figure \ref{fig:prom_targets_after}.

        \begin{figure}[H]
            \centering
            \frame{\includegraphics[width=0.88\columnwidth]{prometheus_targets_green}}
            \caption{Cibles Prometheus}
            \label{fig:prom_targets_after}
        \end{figure}

        \subsubsection{Dashboards Grafana}

        Trois dashboards Grafana \cite{Grafana} ont été provisionnés pour assurer la
        surveillance complète de la plateforme. La figure \ref{fig:grafana_vps} présente
        le dashboard VPS Host Monitoring, basé sur les métriques node-exporter \cite{NodeExporter}.

        \begin{figure}[H]
            \centering
            \frame{\includegraphics[width=0.92\columnwidth]{grafana_vps_host_dashboard}}
            \caption{Dashboard Grafana VPS Host Monitoring}
            \label{fig:grafana_vps}
        \end{figure}

        La figure \ref{fig:grafana_agent} présente le dashboard DevSecOps AI Agent
        dans sa première version opérationnelle (Sprint~5) : trente-trois panneaux
        répartis en cinq lignes couvrant la santé VPS, le pipeline PR, le chat
        opérationnel, les LLM Ollama et les I/O réseau. Des panneaux ont été ajoutés
        aux sprints suivants ; le dashboard comptait quarante-cinq panneaux avant
        l'extension du Sprint~8.1 décrite plus bas.

        \begin{figure}[H]
            \centering
            \frame{\includegraphics[width=0.92\columnwidth]{grafana_agent_dashboard}}
            \caption{Dashboard Grafana DevSecOps AI Agent}
            \label{fig:grafana_agent}
        \end{figure}

        La figure \ref{fig:grafana_pr} présente le dashboard PR Security Reviews, alimenté
        par les données PostgreSQL, avec les panneaux de volume de revues, distribution
        des risques et répartition des verdicts.

        \begin{figure}[H]
            \centering
            \frame{\includegraphics[width=0.88\columnwidth]{grafana_pr_reviews_dashboard}}
            \caption{Dashboard Grafana PR Security Reviews}
            \label{fig:grafana_pr}
        \end{figure}

    \subsection{Sprint 6 — Monitoring host et optimisation du chat}

        \begin{description}
            \itemsep 0.3em
            \item[Objectif du sprint :] Étendre la visibilité aux métriques réelles
                du système hôte et accélérer les réponses de l'assistant de chat.
            \item[Obstacles rencontrés :] \texttt{node-exporter} en réseau Docker
                bridge ne pouvait pas accéder aux systèmes de fichiers \texttt{/proc}
                et \texttt{/sys} du host, le privant des métriques essentielles.
            \item[Livrable :] \texttt{node-exporter} déployé en \texttt{pid:host}
                + \texttt{network\_\allowbreak{}mode:host} avec règle iptables associée, chat
                basculé sur le modèle \texttt{qwen2.\allowbreak{}5-coder:7b} (43\,\% plus rapide
                que le 14B), prompt système compressé de 36\,\% (11\,794 à
                7\,577 caractères).
        \end{description}

        \subsubsection{Déploiement de node-exporter}

        Le déploiement de node-exporter \cite{NodeExporter} a nécessité une configuration
        spécifique car ce
        conteneur doit s'exécuter en dehors du réseau Docker bridge pour avoir accès aux
        métriques réelles du système hôte. Il utilise \texttt{pid: host} et
        \texttt{network\_\allowbreak{}mode: host} pour accéder aux systèmes de fichiers \texttt{/proc}
        et \texttt{/sys} du host. Une règle iptables a été ajoutée pour autoriser le
        bridge Docker à joindre le port 9100 du host. La figure \ref{fig:node_exporter}
        présente la configuration déployée.

        \begin{figure}[H]
            \centering
            \frame{\includegraphics[width=0.88\columnwidth]{node_exporter_config}}
            \caption{Configuration node-exporter}
            \label{fig:node_exporter}
        \end{figure}

        \subsubsection{Optimisation du chat agent}

        C'est au cours de ce sprint qu'a été conduit le benchmark des quatre modèles
        disponibles, dont les résultats ont été présentés au chapitre~2 (tableau
        \ref{tab:llm_benchmark}). À précision égale (80\,\%), le 7b est 43\,\% plus
        rapide par token que le 14b ; il a donc été retenu comme modèle par défaut du
        chat. La compression du prompt système de 11\,794 à 7\,577 caractères
        (36\,\% de moins) a en outre réduit le cache KV et accéléré d'autant les
        réponses.

    \subsection{Sprint 7 — Audit VPS et durcissement}

        \begin{description}
            \itemsep 0.3em
            \item[Objectif du sprint :] Conduire un audit complet de la plateforme
                et éliminer les hallucinations de l'assistant de chat.
            \item[Obstacles rencontrés :] Cinq découvertes critiques. VictoriaMetrics
                était en panne depuis neuf jours, AlertManager n'avait reçu aucune
                alerte depuis le déploiement initial, \texttt{GET /} retournait 404,
                les \textit{live dashboards} Grafana étaient cassés, et le chat UI
                était publiquement accessible sans authentification.
            \item[Livrable :] Cinq correctifs déployés, système anti-hallucination à
                six couches activé sur le chat, authentification HTTP Basic Auth
                déployée sur nginx.
        \end{description}

        \subsubsection{Audit VPS du 28 avril 2026}

        Un audit complet de la plateforme a révélé cinq découvertes critiques requérant
        une correction immédiate. Le tableau \ref{tab:audit} présente ces découvertes et
        les corrections appliquées.

        \begin{table}[H]
        \centering
        \begin{tabular}{|m{3.5cm}|m{4cm}|m{5cm}|}
        \hline
        \textbf{Découverte} & \textbf{Cause racine} & \textbf{Correction} \\
        \hline
        VictoriaMetrics en panne 9 jours & Panic disque-plein, \texttt{restart: unless-stopped} non relancé après libération disque & \texttt{docker compose restart}, 10,9M lignes récupérées \\
        \hline
        AlertManager sans alertes reçues & \texttt{path\_\allowbreak{}prefix} manquant dans \texttt{prometheus.yml} & Ajout \texttt{path\_\allowbreak{}prefix: /alertmanager/} \\
        \hline
        GET / retourne 404 & Absence de location par défaut dans nginx & Ajout \texttt{return 301 /ui} \\
        \hline
        Grafana live dashboards cassés & Headers WebSocket manquants dans nginx & Ajout Upgrade et Connection headers \\
        \hline
        Chat UI publiquement accessible & Absence d'authentification & HTTP Basic Auth via nginx htpasswd \\
        \hline
        \end{tabular}
        \caption{Découvertes et corrections de l'audit VPS du 28 avril 2026}
        \label{tab:audit}
        \end{table}

        La figure \ref{fig:alertmanager_fix} illustre la correction de la configuration
        AlertManager dans \texttt{prometheus.\allowbreak{}yml}. L'absence du paramètre
        \texttt{path\_\allowbreak{}prefix: /\allowbreak{}alertmanager/} entraînait l'envoi de chaque alerte vers
        \texttt{/\allowbreak{}api/\allowbreak{}v2/\allowbreak{}alerts} au lieu de \texttt{/\allowbreak{}alertmanager/\allowbreak{}api/\allowbreak{}v2/\allowbreak{}alerts}, générant
        une erreur 404 silencieuse depuis le premier déploiement.

        \begin{figure}[H]
            \centering
            \frame{\includegraphics[width=0.88\columnwidth]{alertmanager_fix}}
            \caption{Correction AlertManager}
            \label{fig:alertmanager_fix}
        \end{figure}

        \subsubsection{Implémentation du système anti-hallucination}

        Le système anti-hallucination a nettement modifié le comportement du chat.
        Avant les corrections du sprint 7, le modèle inventait des valeurs de métriques
        plausibles sans appeler aucun outil. Après les corrections, toute question portant
        sur une donnée live déclenche systématiquement l'appel d'outil correspondant, et la
        réponse ne cite que des valeurs présentes dans les observations.

        La figure \ref{fig:chat_ui} présente l'interface du BTE Security AI Agent avec
        les blocs de raisonnement et d'appels d'outils visibles ; cette présentation rend
        transparent le processus de génération de la réponse.

        \begin{figure}[H]
            \centering
            \frame{\includegraphics[width=0.92\columnwidth]{chat_ui_screenshot}}
            \caption{Interface BTE Security AI Agent}
            \label{fig:chat_ui}
        \end{figure}

        Au-delà des questions d'état système, l'assistant de chat sait également
        interroger la base de connaissances sécurité hébergée dans PostgreSQL
        grâce à son outil \texttt{query\_\allowbreak{}database}. La figure \ref{fig:db_query_via_chat}
        illustre une requête en langage naturel portant sur l'historique des
        revues de \textit{Pull Requests} : le modèle décide d'invoquer l'outil,
        construit la requête SQL appropriée, exécute la lecture sur la table
        \texttt{pr\_reviews}, puis présente les résultats sous forme de tableau
        markdown. Conformément au système anti-hallucination, chaque valeur
        affichée provient verbatim de l'observation injectée après l'appel
        d'outil, sans qu'aucune donnée ne soit inventée.

        \begin{figure}[H]
            \centering
            \frame{\includegraphics[width=0.92\columnwidth]{db_query_via_chat}}
            \caption{Interrogation de la base de connaissances sécurité via le chat opérationnel}
            \label{fig:db_query_via_chat}
        \end{figure}

        Cette interrogation conversationnelle de la base de connaissances a un
        intérêt opérationnel direct : sans connaître la structure des tables ni
        écrire de SQL, un membre du pôle Sécurité peut consulter l'historique des
        revues, repérer les dépôts les plus exposés ou suivre l'évolution des
        scores de risque dans le temps.

    \subsection{Sprint 8 — Robustesse du parser de revue PR}

        \begin{description}
            \itemsep 0.3em
            \item[Objectif du sprint :] Stabiliser la production des commentaires
                inline GitHub sur l'ensemble des \textit{Pull Requests} réelles.
            \item[Obstacles rencontrés :] Bug critique observé sur PR \#14 et PR
                \#15. Le commentaire global s'affichait correctement mais aucun
                commentaire inline n'apparaissait. L'analyse approfondie a révélé
                trois bugs cumulatifs dans le parser de revue.
            \item[Livrable :] Parser réécrit avec trois stratégies imbriquées
                (\textit{regex} tolérant aux espaces, \textit{walker brace-depth},
                fallback markdown), prompt combiné restructuré pour exiger le bloc
                JSON en première position, bugs reproduits puis vérifiés corrigés
                sur les PRs \#14 et \#15.
        \end{description}

        Lors de la phase de validation finale, deux \textit{Pull Requests} successives
        (PR \#14 et PR \#15) ont produit des revues de sécurité où le commentaire global
        s'affichait correctement, mais où aucun commentaire inline n'apparaissait sur
        l'onglet \textit{Files Changed}. L'analyse approfondie des journaux a révélé trois
        bugs cumulatifs dans le code de parsing du nœud \texttt{analyze\_\allowbreak{}review\_\allowbreak{}node}.

        \subsubsection{Investigation du bug}

        Le journal de l'agent indiquait systématiquement \texttt{comments=0} et un
        \texttt{risk\_score} forcé à la valeur par défaut \texttt{MEDIUM} alors que le
        modèle 14B avait correctement identifié des vulnérabilités de niveau \texttt{HIGH}
        dans son texte. La figure \ref{fig:parser_bug_analysis} illustre les trois
        défauts identifiés et leurs interactions.

        \begin{figure}[H]
            \centering
            \frame{\includegraphics[width=0.88\columnwidth]{parser_bug_analysis}}
            \caption{Analyse des trois bugs cumulatifs du parser PR review (Sprint 8)}
            \label{fig:parser_bug_analysis}
        \end{figure}

        \begin{table}[H]
        \centering
        \begin{tabular}{|m{4cm}|m{4cm}|m{5cm}|}
        \hline
        \textbf{Bug} & \textbf{Symptôme} & \textbf{Cause racine} \\
        \hline
        Regex intolérant aux espaces & \texttt{json\_\allowbreak{}match=None} → défauts \texttt{MEDIUM/\allowbreak{}REQUEST\_\allowbreak{}CHANGES}, \texttt{raw\_\allowbreak{}comments=[]} & Le motif \texttt{\textbackslash\{"risk\_score"...} exigeait que \texttt{"risk\_\allowbreak{}score"} suive immédiatement \texttt{\{}. Lorsque le modèle pré-formattait le JSON (\texttt{\{} + saut de ligne + indentation), le regex échouait silencieusement. \\
        \hline
        \texttt{json.loads} strict sur les caractères en fin & \texttt{JSONDecodeError} → fallback aux groupes regex mais \texttt{raw\_\allowbreak{}comments} restait vide & \texttt{json.\allowbreak{}loads(raw[json\_\allowbreak{}start:])} consommait toute la chaîne jusqu'à la fin du fichier, y compris les délimiteurs Markdown \verb|```| en queue. \\
        \hline
        Modèle omettant le JSON & PR \#15 : la revue se terminait à la section "Recommandations" sans aucun objet JSON & Le modèle, manquant de budget de génération, tronquait sa sortie après la section markdown, omettant complètement le bloc JSON requis par le prompt. \\
        \hline
        \end{tabular}
        \caption{Les trois bugs cumulatifs identifiés sur PR \#14 et PR \#15}
        \label{tab:parser_bugs}
        \end{table}

        \subsubsection{Réécriture du parser et restructuration du prompt}

        La résolution a nécessité une réécriture complète du code de parsing et une
        restructuration du prompt combiné. La figure \ref{fig:parser_rewrite} présente
        les trois stratégies imbriquées du nouveau parser.

        \begin{figure}[H]
            \centering
            \frame{\includegraphics[width=0.40\columnwidth,height=0.85\textheight,keepaspectratio]{parser_rewrite}}
            \caption{Architecture du parser réécrit}
            \label{fig:parser_rewrite}
        \end{figure}

        \begin{lstlisting}[language=Python, caption={Parser PR review réécrit — regex tolérant + walker brace-depth + fallback markdown},
                           label={lst:parser_rewrite}]
# Strategie 1 : regex tolerant aux espaces (correction Sprint 8)
json_match = re.search(
    r'\{\s*"risk_score"\s*:\s*"(?P<rs>CRITICAL|HIGH|MEDIUM|LOW|INFO)"\s*,'
    r'\s*"verdict"\s*:\s*"(?P<v>APPROVE|REQUEST_CHANGES|BLOCK)"',
    raw, re.DOTALL,
)

if json_match:
    risk_score = json_match.group("rs")
    verdict = json_match.group("v")

    # Strategie 2 : walker brace-depth pour trouver l'accolade fermante exacte
    json_start = json_match.start()
    depth, in_string, escape, json_end = 0, False, False, -1
    for i in range(json_start, len(raw)):
        ch = raw[i]
        if escape: escape = False; continue
        if ch == "\\": escape = True; continue
        if ch == '"' and not escape: in_string = not in_string; continue
        if in_string: continue
        if ch == "{": depth += 1
        elif ch == "}":
            depth -= 1
            if depth == 0:
                json_end = i + 1
                break
    if json_end > json_start:
        data = json.loads(raw[json_start:json_end])
        code_review_summary = data.get("code_review_summary", "")
        raw_comments = data.get("comments", [])
else:
    # Strategie 3 : fallback markdown si JSON entierement absent
    md_risk = re.search(
        r'\*?\*?Risk\s*(?:Score)?\s*:\s*\*?\*?\s*'
        r'(CRITICAL|HIGH|MEDIUM|LOW|INFO)', raw, re.IGNORECASE)
    md_verdict = re.search(
        r'\*?\*?Verdict\s*:\s*\*?\*?\s*'
        r'(APPROVE|REQUEST_CHANGES|BLOCK)', raw, re.IGNORECASE)
    if md_risk: risk_score = md_risk.group(1).upper()
    if md_verdict: verdict = md_verdict.group(1).upper()
        \end{lstlisting}

        En complément, le prompt combiné a été restructuré pour exiger la production
        du bloc JSON \textbf{en première position} dans la réponse du modèle, suivi
        du markdown de la revue. Cette inversion empêche le modèle de tronquer la
        sortie avant d'avoir produit les métadonnées critiques. Les figures
        \ref{fig:prompt_restructured_step_1} et \ref{fig:prompt_restructured_step_2}
        présentent respectivement les deux étapes du nouveau prompt STEP~1 / STEP~2.

        \begin{figure}[H]
            \centering
            \frame{\includegraphics[width=0.88\columnwidth]{prompt_restructured_step_1}}
            \caption{Prompt restructuré — STEP 1}
            \label{fig:prompt_restructured_step_1}
        \end{figure}

        \begin{figure}[H]
            \centering
            \frame{\includegraphics[width=0.88\columnwidth]{prompt_restructured_step_2}}
            \caption{Prompt restructuré — STEP 2}
            \label{fig:prompt_restructured_step_2}
        \end{figure}

        \subsubsection{Vérification}

        Le tableau \ref{tab:parser_verification} récapitule les vérifications effectuées
        après la réécriture, validant la robustesse du parser sur les trois cas
        d'échec observés.

        \begin{table}[H]
        \centering
        \begin{tabular}{|m{6cm}|m{6cm}|}
        \hline
        \textbf{Test} & \textbf{Résultat} \\
        \hline
        Regex tolérant aux JSON pré-formatté & \checkmark \texttt{HIGH/BLOCK} extrait correctement \\
        \hline
        Walker tolère les fences \verb|```| en queue & \checkmark Stoppé à la bonne accolade \\
        \hline
        Fallback markdown sur PR \#15 (sans JSON) & \checkmark \texttt{HIGH/BLOCK} extrait du markdown \\
        \hline
        Lignes 12, 17, 26, 34 du diff PR \#14 & \checkmark Toutes valides après parsing \\
        \hline
        Bug identique dans \texttt{\_\allowbreak{}build\_\allowbreak{}degraded\_\allowbreak{}review()} & \checkmark Patché aux deux endroits \\
        \hline
        \end{tabular}
        \caption{Vérifications post-réécriture du parser PR review}
        \label{tab:parser_verification}
        \end{table}

    \subsection{Post-Sprint 8 — Évaluation d'un second backend d'inférence}

        \begin{description}
            \itemsep 0.3em
            \item[Objectif du sprint :] Comparer Ollama à un moteur d'inférence
                alternatif (LocalAI) sur exactement le même fichier GGUF afin de
                transformer le choix initial en mesure objective.
            \item[Obstacles rencontrés :] Bug \textit{cold-load} sur le modèle
                Phi-4. Le \textit{watchdog} \texttt{stream\_\allowbreak{}chunk\_\allowbreak{}timeout=120s}
                de \texttt{langchain\_\allowbreak{}openai} tuait le \textit{stream} avant le
                premier token, car le chargement à froid d'un modèle 15B sur CPU
                dépassait régulièrement deux minutes.
            \item[Livrable :] Sandbox LocalAI opérationnelle en parallèle d'Ollama,
                routeur de chat dual-backend (\texttt{ollama/} ou \texttt{localai/}),
                bug \textit{cold-load} corrigé par désactivation du \textit{watchdog}
                et \textit{pre-warm} explicite du modèle, benchmark croisé livré
                (Ollama +22\,\% de débit sur modèle identique).
        \end{description}

        Ollama \cite{Ollama} avait été choisi en Sprint~1 pour sa réputation
        d'efficacité sur CPU, sans jamais avoir été mesuré contre un moteur
        concurrent servant exactement le même modèle. Ce micro-sprint visait à
        transformer cette intuition en mesure, et accessoirement à vérifier que le
        routeur de chat pouvait accueillir un second backend sans réécriture de la
        boucle ReAct. LocalAI \cite{LocalAI} a été retenu comme point de comparaison
        parce qu'il consomme nativement les fichiers GGUF \cite{GGUF} du Hub Hugging
        Face \cite{HuggingFaceHub}, ce qui donne au passage accès à des modèles
        absents de la bibliothèque Ollama, notamment Phi-4 et Gemma 3. Le travail a
        été mené dans une sandbox isolée, sans aucune modification du stack de
        production : le pipeline de revue de PR continue de tourner exclusivement sur
        Ollama, et LocalAI ne démarre qu'à la demande via un fichier Compose
        distinct.

        \subsubsection{Mise en place de la sandbox LocalAI}

        La sandbox est définie dans le fichier \texttt{docker-compose.\allowbreak{}localai.\allowbreak{}yml},
        séparé du Compose principal, qui rejoint le réseau Docker existant et expose
        une API compatible OpenAI sur le port hôte 8081 (douze threads, contexte de
        8\,192 tokens, limite mémoire de 24~Go). Les modèles de la galerie LocalAI
        s'installent par un simple \texttt{POST /\allowbreak{}models/\allowbreak{}apply} ; les GGUF hors
        galerie, comme le \texttt{qwen2.\allowbreak{}5-coder-7b} choisi pour répliquer le modèle
        servi par Ollama, demandent un fichier YAML décrivant la source et les
        paramètres d'exécution.

        \begin{lstlisting}[style=yaml, caption={Configuration LocalAI du modèle Qwen2.5-Coder 7B, identique à celui servi par Ollama}, label={lst:localai_qwen_yaml}]
name: qwen2.5-coder-7b
parameters:
  model: "huggingface://bartowski/Qwen2.5-Coder-7B-Instruct-GGUF/\
Qwen2.5-Coder-7B-Instruct-Q4_K_M.gguf"
context_size: 8192
threads: 12
f16: true
mmap: true
        \end{lstlisting}

        Le format ChatML de la famille Qwen 2.5 est détecté automatiquement par
        LocalAI à partir des métadonnées du fichier GGUF ; la cohérence des réponses
        produites par les deux backends sur un même prompt a confirmé après coup que
        cette détection fonctionnait.

        \subsubsection{Architecture du routeur dual-backend}

        Le routeur de chat (fichier \texttt{agent/\allowbreak{}app/\allowbreak{}routers/\allowbreak{}chat.\allowbreak{}py}) a été étendu
        pour accepter un identifiant de modèle préfixé par le backend cible :
        \texttt{ollama/\allowbreak{}qwen2.\allowbreak{}5-coder:7b} continue d'emprunter le chemin de production,
        tandis que \texttt{localai/\allowbreak{}phi-4} ou \texttt{localai/\allowbreak{}qwen2.\allowbreak{}5-coder-7b}
        bifurque vers la sandbox. En l'absence de préfixe, le routeur retombe sur
        Ollama, ce qui préserve les conversations sauvegardées avant l'ajout de
        LocalAI. La route \texttt{/\allowbreak{}chat/\allowbreak{}models} fusionne les deux catalogues dans le
        même menu de l'interface ; si la sandbox est arrêtée, seuls les modèles
        Ollama apparaissent, de sorte que la défaillance du backend optionnel
        n'affecte jamais celui de production.

        \subsubsection{Découverte et résolution d'un bug de \og cold-load \fg{}}

        Le premier essai réel a échoué : après sélection du modèle
        \texttt{localai/\allowbreak{}phi-4}, le flux SSE s'interrompait après exactement
        cent vingt secondes. Le client \texttt{ChatOpenAI} de la bibliothèque
        \texttt{langchain\_\allowbreak{}openai} embarque en effet un chien de garde par chunk de
        cent vingt secondes, indépendant du timeout global de la requête. Or, sur un
        CPU sans GPU, le chargement à froid d'un modèle 15B quantisé comme Phi-4
        (8,5 Go sur disque), suivi du traitement d'un prompt système d'environ
        1\,900 tokens, dépasse régulièrement ce délai avant que le premier token ne
        soit produit : le chien de garde tuait le flux pendant que le modèle
        finissait de se charger en RAM.

        La correction combine trois ajustements. Le chien de garde a été désactivé
        (\texttt{stream\_\allowbreak{}chunk\_\allowbreak{}timeout=None}) : entre deux conteneurs du même
        \textit{bridge} Docker, le \textit{keepalive} TCP du noyau suffit à détecter
        une connexion morte. Un pré-chargement explicite du modèle a été ajouté avant
        le démarrage du streaming, sous la forme d'une petite requête synchrone qui
        force LocalAI à amener le GGUF en mémoire. Enfin, deux événements SSE
        intermédiaires (\og Warming up\ldots{} \fg{} puis \og Generating\ldots{} \fg{})
        renseignent l'utilisateur pendant le chargement à froid, durant lequel
        l'interface restait auparavant figée.

        \subsubsection{Benchmark croisé sur modèle identique}

        Une fois le routeur stabilisé, nous avons construit un protocole de mesure
        pour comparer les deux backends. Comparer un modèle dense Qwen 2.5 servi par
        Ollama contre un MoE Qwen 3 servi par LocalAI aurait mélangé deux écarts
        distincts, celui des architectures de modèles et celui des moteurs
        d'inférence. Nous voulions isoler le second. Le fichier YAML LocalAI pointe
        donc sur exactement le même GGUF que celui distribué par Ollama, à savoir
        \texttt{bartowski/\allowbreak{}Qwen2.\allowbreak{}5-Coder-7B-Instruct-Q4\_\allowbreak{}K\_\allowbreak{}M.\allowbreak{}gguf}, soit 4,7 Go.

        Le script \texttt{scripts/\allowbreak{}benchmark-backends.\allowbreak{}sh} reprend la méthodologie du
        banc d'essai mono-backend introduit en Sprint 6. Une passe à froid de quatre
        tokens, non chronométrée, amène d'abord le modèle en RAM. La passe chaude
        qui suit, sur un prompt identique demandant un résumé OWASP Top 10 d'une
        application PHP à température 0,1, est mesurée au chronomètre et limitée à
        quatre-vingts tokens. Le débit est ensuite calculé en parsant le champ
        \texttt{usage.\allowbreak{}completion\_\allowbreak{}tokens} de la réponse JSON. Pour rendre la
        comparaison équitable, les deux appels sont émis depuis l'intérieur du
        conteneur \texttt{devsecops-agent} avec la bibliothèque \texttt{httpx}.
        Ollama n'étant pas publié sur le port hôte et son image officielle
        n'embarquant pas \texttt{curl}, cette approche commune élimine tout biais
        de surcoût d'environnement entre les deux mesures.

        Les résultats obtenus le 12 mai 2026 sur le VPS de production sont
        présentés dans le tableau \ref{tab:benchmark_backends}.

        \begin{table}[H]
        \centering
        \begin{tabular}{|m{2.5cm}|m{4.5cm}|c|c|c|}
        \hline
        \textbf{Backend} & \textbf{Modèle} & \textbf{Tokens} & \textbf{Temps (s)} & \textbf{Débit (tok/s)} \\
        \hline
        Ollama  & \texttt{qwen2.\allowbreak{}5-coder:7b}  & 69 & 12,56 & \textbf{5,49} \\
        \hline
        LocalAI & \texttt{qwen2.\allowbreak{}5-coder-7b}  & 80 & 17,79 & \textbf{4,50} \\
        \hline
        \end{tabular}
        \caption{Benchmark croisé Ollama vs LocalAI sur modèle identique (Qwen2.5-Coder 7B, quantification Q4\_K\_M, prompt OWASP Top 10, douze cœurs Haswell, sans GPU)}
        \label{tab:benchmark_backends}
        \end{table}

        \subsubsection{Interprétation et décision}

        Sur le même matériel et avec exactement le même fichier de poids, Ollama
        produit un débit supérieur de 22\,\% à celui de LocalAI (5,49 contre 4,50
        tokens par seconde). Les deux backends produisent des énumérations OWASP
        sémantiquement équivalentes. La légère différence de longueur vient du fait
        qu'Ollama s'arrête naturellement après les dix risques, tandis que LocalAI
        continue d'élaborer jusqu'à atteindre la limite des quatre-vingts tokens.
        L'écart de débit s'explique vraisemblablement par l'intégration plus serrée
        d'Ollama avec \texttt{llama.cpp} et par les options qu'il active à
        l'exécution (\texttt{OLLAMA\_\allowbreak{}FLASH\_\allowbreak{}ATTENTION=1},
        \texttt{OLLAMA\_\allowbreak{}KV\_\allowbreak{}CACHE\_\allowbreak{}TYPE=q8\_\allowbreak{}0}, \texttt{OLLAMA\_\allowbreak{}NUM\_\allowbreak{}THREAD=12}),
        contre une couche d'orchestration plus générique côté LocalAI qui doit
        couvrir plusieurs runtimes.

        Ce résultat ne modifie pas le choix opérationnel : Ollama reste le backend de
        production. Il le confirme en revanche sur une base mesurée, alors qu'il reposait
        jusqu'ici sur la réputation de l'outil et sur notre habitude de l'employer. LocalAI
        conserve une place utile comme sandbox optionnelle, pour accéder à des modèles
        HuggingFace absents de la bibliothèque Ollama et pour préparer une comparaison
        future, le jour où un nouveau moteur d'inférence deviendra pertinent.

    \subsection{Sprint 8.1 — Expansion de l'observabilité et correction du \textit{template} Phi-4}

        \begin{description}
            \itemsep 0.3em
            \item[Objectif du sprint :] Étendre la couverture d'observabilité à
                chaque conteneur du VPS et instrumenter le routeur de chat pour
                mesurer en production les performances comparées des deux backends.
            \item[Obstacles rencontrés :] cAdvisor architecturalement incompatible
                avec le \textit{containerd snapshotter} du Docker 29.4
                (\texttt{failed to identify the read-write layer ID}) ; bug latent
                du modèle Phi-4 retournant silencieusement zéro token à chaque
                requête (\textit{template} ChatML mal formaté).
            \item[Livrable :] \textit{Poller} Docker autonome interrogeant le
                socket toutes les trente secondes (six métriques par conteneur),
                instrumentation A/B du routeur de chat (quatre métriques
                \texttt{chat\_*} labellisées par couple \texttt{backend} et
                \texttt{model}), cinq nouvelles \textit{rows} ajoutées au dashboard
                Grafana de l'agent, \textit{template} Phi-4 corrigé avec injection
                explicite du \textit{prompt} de génération.
        \end{description}

        Le micro-sprint précédent a laissé une question ouverte : comment suivre en
        production, et non plus seulement au banc d'essai, le comportement de chaque
        modèle de chat ? La couverture d'observabilité héritée des sprints précédents
        s'arrêtait à l'agent et à Ollama, sans visibilité par conteneur ni comparatif
        entre les deux backends. Le sprint 8.1 a donc instrumenté chaque conteneur du
        VPS, ajouté une télémétrie A/B au routeur de chat, et mis à profit cette
        instrumentation pour corriger un bug latent du modèle Phi-4.

        \subsubsection{Tentative de déploiement de cAdvisor et pivot architectural}

        La voie naturelle pour obtenir des métriques par conteneur était
        \textit{cAdvisor}, le projet de référence de Google pour cette collecte. Les
        essais menés sur deux versions et trois combinaisons de \textit{flags} ont
        abouti à un constat sans appel : le démon Docker 29.4 utilise le
        \textit{containerd snapshotter} comme magasin d'images, qui ne maintient pas
        la base de couches dont cAdvisor a besoin, et chaque \textit{cgroup}
        découvert était aussitôt écarté avec l'erreur \texttt{failed to identify the
        read-write layer ID}. Plutôt que d'insister, nous avons exploité une
        ressource déjà disponible : l'agent monte le socket Docker pour ses outils
        d'introspection. Une centaine de lignes Python suffisent à interroger l'API
        Docker Engine toutes les trente secondes et à exposer le résultat en jauges
        Prometheus, sans conteneur supplémentaire, indépendamment du pilote de
        stockage et sans privilège nouveau.

        \begin{lstlisting}[language=Python, caption={Cœur du \textit{poller} docker-stats — calcul CPU\% relatif à un cœur}, label={lst:docker_stats_poller}]
async def _poll_docker_stats():
    transport = httpx.AsyncHTTPTransport(uds="/var/run/docker.sock")
    while True:
        async with httpx.AsyncClient(transport=transport,
                                     base_url="http://docker",
                                     timeout=10.0) as client:
            for c in (await client.get("/containers/json")).json():
                cid, name = c["Id"], c["Names"][0].lstrip("/")
                img = c.get("Image", "?")
                st = (await client.get(f"/containers/{cid}/stats",
                                       params={"stream": "false"})).json()
                # Memoire
                mem = st["memory_stats"]
                container_memory_bytes.labels(
                    name=name, image=img).set(mem.get("usage", 0))
                # CPU% (rapport delta-CPU / delta-systeme x cores)
                cpu, pre = st["cpu_stats"], st["precpu_stats"]
                d_cpu = (cpu["cpu_usage"]["total_usage"]
                         - pre["cpu_usage"]["total_usage"])
                d_sys = (cpu.get("system_cpu_usage", 0)
                         - pre.get("system_cpu_usage", 0))
                online = cpu.get("online_cpus") or 1
                pct = (d_cpu / d_sys) * online * 100.0 if d_sys > 0 else 0.0
                container_cpu_percent.labels(
                    name=name, image=img).set(max(0.0, pct))
        await asyncio.sleep(30)
        \end{lstlisting}

        Le premier passage du \textit{poller} a immédiatement révélé douze conteneurs
        actifs, avec une consommation mémoire conforme à l'intuition : l'agent occupe
        environ 127 Mo, le conteneur LocalAI grimpait jusqu'à 25,7 Go lorsque le modèle
        Phi-4 était chargé, et Ollama tournait autour de 5,8 Go en veille. Cette
        granularité, indisponible jusqu'alors, permet désormais de détecter une dérive de
        mémoire ou une consommation CPU anormale sur n'importe lequel des services de la
        plateforme.

        \subsubsection{Instrumentation du routeur de chat pour la télémétrie A/B}

        La fonction \texttt{\_\allowbreak{}stream\_\allowbreak{}react()} du routeur de chat a été enveloppée d'un
        bloc \texttt{try/\allowbreak{}except/\allowbreak{}finally} qui mesure trois grandeurs pour chaque appel
        à \texttt{/\allowbreak{}chat/\allowbreak{}stream} : la latence avant le premier \textit{token} émis
        (\textit{Time-To-First-Token}), la durée totale du \textit{stream}, et le
        nombre de \textit{chunks} effectivement renvoyés à l'interface. Chaque
        observation est labellisée par le couple \texttt{(backend, model)}, de sorte
        que Grafana peut tracer côte à côte la performance d'Ollama et de LocalAI sur
        la même question. L'instrumentation est strictement passive : aucune branche
        logique de la boucle ReAct n'a été modifiée, et le chemin de revue de Pull
        Request, qui n'emprunte pas \texttt{\_\allowbreak{}stream\_\allowbreak{}react()}, reste inchangé.

        \subsubsection{Cinq nouvelles \textit{rows} sur le dashboard Grafana}

        Le Sprint~8.1 ajoute cinq \textit{rows} au dashboard de l'agent (soixante-six
        panneaux au total, contre trente-trois lors du Sprint~5), ainsi qu'un
        indicateur LocalAI UP/DOWN dans la bande d'en-tête, à côté de l'indicateur
        Ollama existant. Le tableau \ref{tab:grafana_rows_8_1} récapitule
        l'organisation visuelle.

        \begin{table}[H]
        \centering
        \footnotesize
        \begin{tabular}{|p{4.2cm}|p{10.5cm}|}
        \hline
        \textbf{Row} & \textbf{Contenu} \\
        \hline
        LocalAI Sandbox & UP/DOWN \;|\; nombre de modèles installés \;|\; latence du \textit{health probe} \;|\; table des modèles avec taille en Go \\
        \hline
        Chat Backend A/B & Requêtes/min par backend \;|\; \textit{tokens}/s par couple backend×modèle \;|\; TTFT p50/p95 \;|\; durée totale p50/p95 \;|\; table des statuts sur 24\,h \\
        \hline
        Containers & Mémoire par conteneur (Go, stacked) \;|\; CPU \% par conteneur \;|\; trafic RX/TX (octets/s) \;|\; \textit{snapshot} tabulaire \\
        \hline
        Service Health & Onze indicateurs UP/DOWN couvrant agent, ollama, localai, postgres, redis, nginx, prometheus, alertmanager, grafana, victoriametrics, node-exporter \\
        \hline
        Firing Alerts & Compteur total \;|\; table des alertes actives (\texttt{ALERTS\{alertstate="firing"\}}) \\
        \hline
        \end{tabular}
        \caption{Organisation des cinq nouvelles \textit{rows} ajoutées au \textit{dashboard} Grafana lors du sprint 8.1}
        \label{tab:grafana_rows_8_1}
        \end{table}

        Les onze indicateurs de la \textit{row} \og Service Health \fg{} agrègent
        trois sources de signal : la métrique \texttt{up\{job=...\}} de Prometheus
        pour les cibles scrutées directement, les jauges \texttt{ollama\_\allowbreak{}reachable}
        et \texttt{localai\_\allowbreak{}reachable} émises par leurs \textit{pollers}, et la jauge
        \texttt{container\_\allowbreak{}running} du nouveau \textit{poller} Docker pour les
        services non instrumentés Prometheus. L'opérateur dispose ainsi d'un tableau
        vert/rouge immédiat à l'ouverture du \textit{dashboard}.

        \subsubsection{Découverte et correction du bug du \textit{template} Phi-4}

        L'instrumentation A/B a immédiatement révélé une anomalie sur le modèle
        Phi-4 : la requête se terminait avec un statut \texttt{ok}, mais le compteur
        \texttt{chat\_\allowbreak{}tokens\_\allowbreak{}streamed\_\allowbreak{}total\{model="phi-4"\}} restait obstinément à
        zéro. Sans la nouvelle télémétrie, ce comportement aurait été indétectable
        au niveau opérationnel. L'utilisateur voyait simplement une réponse vide,
        sans qu'aucun log ne rapporte d'erreur. La cause racine, localisée dans le
        fichier \texttt{localai/\allowbreak{}config/\allowbreak{}phi-4.\allowbreak{}yaml}, était subtile : la chaîne de
        \textit{template} utilisée pour la clé \texttt{chat} était identique à celle
        de \texttt{chat\_\allowbreak{}message}, ce qui revient à n'enrouler aucun préfixe de
        génération autour de l'historique formaté. Le modèle recevait donc une
        conversation déjà close (terminée par son propre \texttt{<|im\_end|>}) et
        n'avait aucune raison de prendre la parole.

        Le \textit{template} a été récrit en s'alignant strictement sur le fichier
        \texttt{tokenizer\_\allowbreak{}config.\allowbreak{}json} officiel publié par Microsoft sur le Hub
        Hugging Face. La clé \texttt{chat\_\allowbreak{}message} sert désormais uniquement à
        formater un message individuel, tandis que \texttt{chat} injecte l'historique
        complet et ajoute explicitement le \textit{prompt} de génération
        \texttt{<|im\_\allowbreak{}start|>assistant<|im\_\allowbreak{}sep|>}. Les retours à la ligne parasites
        autour de \texttt{<|im\_sep|>} ont également été supprimés, et le suffixe
        \texttt{|-} de YAML strip-newlines a été employé pour garantir qu'aucun
        caractère blanc ne sépare le \textit{prompt} de génération de la première
        sortie du modèle.

        \begin{lstlisting}[style=yaml, caption={Correction du \textit{template} ChatML-with-im\_sep de Phi-4}, label={lst:phi4_fix}]
template:
  chat_message: |-
    <|im_start|>{{ .RoleName }}<|im_sep|>{{ .Content }}<|im_end|>
  chat: |-
    {{.Input}}<|im_start|>assistant<|im_sep|>
  completion: |
    {{.Input}}
        \end{lstlisting}

        Le fichier corrigé est projeté dans le conteneur par un \textit{bind-mount}
        déclaré dans \texttt{docker-compose.\allowbreak{}localai.\allowbreak{}yml}, ce qui le fait survivre
        à un cycle complet de \texttt{docker compose down/\allowbreak{}up}. La validation a porté
        sur les deux chemins : en appel direct à LocalAI, \texttt{completion\_\allowbreak{}tokens}
        vaut désormais 3 là où il valait 0 ; à travers le routeur de chat de l'agent,
        le compteur \texttt{chat\_\allowbreak{}tokens\_\allowbreak{}streamed\_\allowbreak{}total\{model="phi-4"\}} est passé
        de 0 à 51 \textit{tokens}.

        \subsubsection{Synthèse de la couverture observabilité}

        Le tableau \ref{tab:metrics_8_1_summary} récapitule la couverture de
        monitoring obtenue à l'issue du sprint 8.1, en distinguant les métriques
        héritées des sprints antérieurs des quinze nouvelles.

        \begin{table}[H]
        \centering
        \footnotesize
        \begin{tabular}{|p{3.8cm}|p{2cm}|p{8.5cm}|}
        \hline
        \textbf{Famille de métriques} & \textbf{Nombre} & \textbf{Périmètre} \\
        \hline
        Héritées (Sprints 5--8) & 13 & Compteurs et histogrammes de l'agent, jauges Ollama (5), disque (2) \\
        \hline
        \texttt{container\_*} (Sprint~8.1) & 6 & Mémoire, limite mémoire, CPU\%, RX/TX réseau, état UP/DOWN par conteneur \\
        \hline
        \texttt{localai\_*} (Sprint~8.1) & 5 & Joignabilité, nombre de modèles, présence par modèle, taille en Go, latence \textit{health probe} \\
        \hline
        \texttt{chat\_*} (Sprint~8.1) & 4 & Total des requêtes par statut, durée totale, TTFT, \textit{tokens} émis (tous par backend×modèle) \\
        \hline
        \textbf{Total} & \textbf{28} & 13 historiques + 15 nouvelles \\
        \hline
        \end{tabular}
        \caption{Inventaire consolidé des métriques Prometheus à l'issue du sprint 8.1}
        \label{tab:metrics_8_1_summary}
        \end{table}

% ─────────────────────────────────────────────────────────────────────────────
\section{Tests de validation}

    Afin de vérifier le bon fonctionnement de chaque mécanisme de sécurité et de résilience
    de la solution, nous avons conduit une série de tests unitaires sur les composants
    critiques de la plateforme. Ces tests reproduisent des conditions réelles d'erreur ou
    d'attaque pour valider les comportements attendus.

    \subsection{Test de la validation HMAC-SHA256}

        Le premier test vérifie que l'agent rejette correctement les webhooks dont la
        signature est invalide ou absente. La figure \ref{fig:test_hmac} présente le
        résultat du test : un webhook avec une signature incorrecte reçoit un code HTTP
        403 Forbidden, tandis qu'un webhook correctement signé est accepté avec un code 202.

        \begin{figure}[H]
            \centering
            \frame{\includegraphics[width=0.88\columnwidth]{test_hmac_validation}}
            \caption{Test HMAC}
            \label{fig:test_hmac}
        \end{figure}

    \subsection{Test de la déduplication Redis}

        Le deuxième test vérifie que deux webhooks identiques (même dépôt, même numéro
        de PR, même SHA de commit) ne déclenchent qu'un seul pipeline. La figure
        \ref{fig:test_dedup} montre que le second webhook retourne un code 202 mais que
        l'agent l'ignore silencieusement grâce à la clé de déduplication Redis
        \texttt{SET NX} ; le premier pipeline se termine alors sans interruption.

        \begin{figure}[H]
            \centering
            \frame{\includegraphics[width=0.88\columnwidth]{test_deduplication}}
            \caption{Test déduplication Redis}
            \label{fig:test_dedup}
        \end{figure}

    \subsection{Test du circuit breaker LLM}

        Le troisième test vérifie le comportement de l'agent lorsque Ollama est
        indisponible. Après trois échecs consécutifs d'appels LLM, le circuit breaker
        passe en état OPEN et active le fallback de classification regex. La figure
        \ref{fig:test_circuit} montre que le pipeline continue de fonctionner avec une
        classification basée sur les extensions de fichiers, sans bloquer sur
        l'indisponibilité du modèle.

        \begin{figure}[H]
            \centering
            \frame{\includegraphics[width=0.88\columnwidth]{test_circuit_breaker}}
            \caption{Test circuit breaker}
            \label{fig:test_circuit}
        \end{figure}

    \subsection{Test du gardien disque}

        Le quatrième test vérifie le déclenchement automatique du nettoyage Docker
        lorsque l'utilisation du disque dépasse 90\%. En simulant un seuil critique via
        la jauge Prometheus \texttt{agent\_\allowbreak{}disk\_\allowbreak{}used\_\allowbreak{}percent}, nous avons vérifié que
        la commande \texttt{docker builder prune -f} est exécutée automatiquement et
        que le résultat est transmis sur Slack, comme le montre la figure
        \ref{fig:test_disk_guard}.

        \begin{figure}[H]
            \centering
            \frame{\includegraphics[width=0.75\columnwidth]{test_disk_guard}}
            \caption{Test gardien disque}
            \label{fig:test_disk_guard}
        \end{figure}

% ─────────────────────────────────────────────────────────────────────────────
\section{Résultats}

    \subsection{Résultat de la revue sur une Pull Request de démonstration}

        Le pipeline a été validé de bout en bout sur la PR~\#18 du dépôt
        \texttt{GhaiethFerchichi/\allowbreak{}Vunl-application}, ouverte le 28~mai~2026 pour
        un module bancaire de démonstration contenant cinq vulnérabilités types.
        Les figures~\ref{fig:pr_full_result}, \ref{fig:pr_full_result_2}
        et~\ref{fig:pr_full_result_3} en donnent la vue d'ensemble, scindée en
        trois parties pour la lisibilité : commentaires de revue, suggestions
        \textit{inline}, verdict \texttt{BLOCK} et échec du check
        \texttt{security-ai-agent}.

        \begin{figure}[H]
            \centering
            \frame{\includegraphics[trim=0 5109 0 0, clip, width=0.82\columnwidth]{pr_review_full_result}}
            \caption{PR~\#18 — vue d'ensemble de la revue sur GitHub (1/3) :
                     en-tête de la \textit{Pull Request} et premiers commentaires de revue}
            \label{fig:pr_full_result}
        \end{figure}

        \begin{figure}[H]
            \centering
            \frame{\includegraphics[trim=0 2554 0 2555, clip, width=0.82\columnwidth]{pr_review_full_result}}
            \caption{PR~\#18 — vue d'ensemble de la revue sur GitHub (2/3) :
                     suite des commentaires \textit{inline} sur le diff}
            \label{fig:pr_full_result_2}
        \end{figure}

        \begin{figure}[H]
            \centering
            \frame{\includegraphics[trim=0 0 0 5109, clip, width=0.82\columnwidth]{pr_review_full_result}}
            \caption{PR~\#18 — vue d'ensemble de la revue sur GitHub (3/3) :
                     revue de sécurité consolidée et verdict \texttt{BLOCK}}
            \label{fig:pr_full_result_3}
        \end{figure}

        Le tableau~\ref{tab:pr_results} synthétise le résultat de cette revue.

        \begin{table}[H]
        \centering
        \begin{tabular}{|c|c|c|c|c|c|}
        \hline
        \textbf{PR} & \textbf{Risque} & \textbf{Verdict} & \textbf{Classification} & \textbf{Durée} & \textbf{Commentaires} \\
        \hline
        \#18 & HIGH & BLOCK & feature & \textasciitilde15 min & 5 inline \\
        \hline
        \end{tabular}
        \caption[Résultat de la revue PR de démonstration]{Résultat de la revue PR de démonstration — dépôt GhaiethFerchichi/Vunl-application}
        \label{tab:pr_results}
        \end{table}

        Le fichier de démonstration contenait cinq vulnérabilités volontairement
        introduites, ce qui fournit une vérité terrain explicite. Le tableau
        \ref{tab:pr18_detection} confronte ces cinq failles aux \textit{findings}
        effectivement produits par l'agent.

        \begin{table}[H]
        \centering
        \begin{tabular}{|m{5.2cm}|c|c|m{3.2cm}|}
        \hline
        \textbf{Vulnérabilité plantée} & \textbf{Ligne} & \textbf{Détectée} & \textbf{Source} \\
        \hline
        Injection SQL & 6 & \checkmark & Analyse LLM \\
        \hline
        Secrets en dur (clés AWS, mot de passe BD) & 11--13 & \checkmark & Analyse LLM \\
        \hline
        Path traversal & 17 & \checkmark & LLM + Semgrep \\
        \hline
        Injection de commande & 22 & \checkmark & LLM + Semgrep \\
        \hline
        Hachage faible (MD5) & 27 & \checkmark & Analyse LLM \\
        \hline
        \end{tabular}
        \caption{Détection sur la PR~\#18 : cinq vulnérabilités plantées, cinq détectées}
        \label{tab:pr18_detection}
        \end{table}

        Sur ce jeu contrôlé, l'agent a détecté les cinq vulnérabilités plantées (rappel
        de 5/5), sans \textit{finding} erroné dans le fichier modifié. Deux observations
        en découlent. D'abord, les secrets codés
        en dur (lignes 11 à 13) ont été repérés par le modèle alors que Gitleaks ne les
        avait pas signalés : le LLM et les scanners spécialisés se complètent plutôt
        qu'ils ne se recouvrent. Ensuite, ces chiffres portent sur des vulnérabilités
        introduites à dessein et de facture classique ; ils valident la chaîne de
        détection de bout en bout, sans préjuger du comportement de l'agent face à des
        failles plus subtiles ou inédites.

    \subsection{Revues produites sur un corpus de Pull Requests}

        La base de connaissances a accumulé huit revues réelles au fil des sprints,
        toutes conservées dans la table \texttt{pr\_\allowbreak{}reviews}. La figure
        \ref{fig:pr_reviews_corpus} en présente les cinq revues de PRs applicatives
        (classification \textit{feature}), celles pour lesquelles le moteur Semgrep entre
        en jeu : pour chacune, le score de risque, le verdict rendu et le nombre de
        \textit{findings} remontés par Semgrep.

        \begin{figure}[H]
            \centering
            \frame{\includegraphics[width=\textwidth]{pr_reviews_corpus}}
            \caption{Revues de PRs applicatives — risque, verdict et \textit{findings} SAST (table \texttt{pr\_reviews})}
            \label{fig:pr_reviews_corpus}
        \end{figure}

        Semgrep a relevé entre 28 et 37 \textit{findings} de niveau ERROR par PR, et le
        verdict suit le risque : \texttt{REQUEST\_\allowbreak{}CHANGES} pour les deux PRs classées
        \textsc{medium}, \texttt{BLOCK} pour les trois classées \textsc{high}. La PR~\#18,
        détaillée plus haut, en est le cas le plus net : cinq vulnérabilités plantées,
        toutes détectées, verdict \texttt{BLOCK}. Sur CPU, l'inférence du modèle 14B (environ
        3 tokens par seconde) reste le principal facteur de durée d'une revue ; c'est là que
        la migration GPU évoquée en conclusion changerait l'échelle.

    \subsection{Comparaison avec la revue manuelle}

        Le tableau \ref{tab:baseline} met en regard le processus manuel décrit au
        chapitre~1 et la chaîne automatisée mise en place.

        \begin{table}[H]
        \centering
        \begin{tabular}{|m{3.2cm}|m{4.7cm}|m{4.7cm}|}
        \hline
        \textbf{Critère} & \textbf{Revue manuelle (AS-IS)} & \textbf{BTE Security AI Agent (TO-BE)} \\
        \hline
        Délai de revue & Jusqu'à 24~heures, selon la disponibilité et l'expertise du relecteur & 15 à 25 minutes selon la PR \\
        \hline
        Couverture & Variable selon la disponibilité et l'expertise du relecteur & Cinq scanners et analyse LLM, OWASP Top 10 systématique \\
        \hline
        Gate CI/CD & Aucun blocage technique & Verdict \texttt{APPROVE} / \texttt{REQUEST\_\allowbreak{}CHANGES} / \texttt{BLOCK} \\
        \hline
        Secrets et dépendances & Souvent manqués à l'œil nu & Gitleaks et OSV-Scanner systématiques \\
        \hline
        Disponibilité & Tributaire du relecteur du jour & Automatique, à chaque \textit{Pull Request} \\
        \hline
        Traçabilité & Fil de discussion GitHub & Base PostgreSQL et tableaux de bord \\
        \hline
        \end{tabular}
        \caption{Revue de sécurité manuelle (AS-IS) vs automatisée (TO-BE)}
        \label{tab:baseline}
        \end{table}

        L'automatisation ne remplace pas le relecteur humain : elle lui retire le travail
        répétitif de détection systématique et le recentre sur les décisions
        d'architecture et sur les cas que l'agent remonte en \texttt{BLOCK}.

    \subsection{Résultats de l'analyse de conformité et observabilité}

        La figure \ref{fig:prom_alerts_ok} présente l'état des quinze règles d'alerte
        Prometheus, toutes inactives, confirmant l'absence d'incident en production lors
        de la validation finale.

        \begin{figure}[H]
            \centering
            \frame{\includegraphics[width=0.88\columnwidth]{prometheus_alerts_inactive}}
            \caption{Prometheus — 15 règles d'alerte inactives}
            \label{fig:prom_alerts_ok}
        \end{figure}

    \subsection{Validation de la plateforme}

        Le tableau \ref{tab:validation} présente la checklist de validation complète
        effectuée sur la plateforme avant la présentation finale.

        \begin{table}[H]
        \centering
        \begin{tabular}{|m{6cm}|m{5cm}|c|}
        \hline
        \textbf{Vérification} & \textbf{Méthode} & \textbf{Statut} \\
        \hline
        Pipeline PR end-to-end & PR réelle → commentaire en 15–25 min & \checkmark \\
        \hline
        5 scanners exécutés & Artefacts dans ./artifacts/scans/ & \checkmark \\
        \hline
        LLM local (aucune fuite) & Port 11434 interne uniquement & \checkmark \\
        \hline
        4 cibles Prometheus vertes & /prometheus/targets & \checkmark \\
        \hline
        15 règles d'alerte inactives & /prometheus/alerts & \checkmark \\
        \hline
        AlertManager → Slack fonctionnel & Test alert → message Slack reçu & \checkmark \\
        \hline
        3 dashboards Grafana chargés & Données réelles dans chaque panneau & \checkmark \\
        \hline
        VictoriaMetrics données 90j & Requête Grafana > 30 jours & \checkmark \\
        \hline
        Anti-hallucination chat validé & "CPU usage?" → appel outil + valeur réelle & \checkmark \\
        \hline
        Gardien disque actif & Jauge agent\_disk\_used\_percent Prometheus & \checkmark \\
        \hline
        Health digest Slack quotidien & Message reçu à 09h00 UTC & \checkmark \\
        \hline
        Métriques host node-exporter & PromQL node\_memory\_MemAvailable & \checkmark \\
        \hline
        nginx Basic Auth /ui protégé & Accès refusé sans credentials & \checkmark \\
        \hline
        Base de données PR reviews & SELECT pr\_reviews — 8 revues (PR~\#11 à \#18) & \checkmark \\
        \hline
        \end{tabular}
        \caption{Checklist de validation complète de la plateforme}
        \label{tab:validation}
        \end{table}

    \subsection{Limites identifiées et plan de durcissement}

        La checklist précédente couvre le périmètre fonctionnel et observabilité, mais
        la plateforme reste configurée pour l'itération quotidienne plutôt que pour une
        exploitation durcie. Plusieurs services internes ont conservé un \textit{binding}
        sur l'hôte hérité de la phase de mise au point. Ce choix a facilité le diagnostic
        au cours des sprints, dans la mesure où une commande \texttt{curl
        http:/\allowbreak{}/\allowbreak{}vps:9090/\allowbreak{}api/\allowbreak{}v1/\allowbreak{}query} était plus rapide à taper que l'ouverture d'un
        tunnel SSH. Il expose toutefois une surface qu'il faut réduire avant tout
        passage en production réelle. Le tableau \ref{tab:hardening_gaps}
        récapitule les écarts relevés et la correction prévue.

        \begin{table}[H]
        \centering
        \begin{tabular}{|m{2.7cm}|m{2cm}|m{4cm}|m{4cm}|}
        \hline
        \textbf{Service} & \textbf{Port} & \textbf{Risque associé} & \textbf{Correction prévue} \\
        \hline
        Redis & 6379 & Lecture du cache, des clés de déduplication et du rate limiter & Retrait du binding hôte, mot de passe \texttt{--requirepass} \\
        \hline
        Prometheus & 9090 & Lecture des métriques, des cibles et des règles d'alerte & Retrait du binding hôte, accès passé exclusivement par nginx \\
        \hline
        AlertManager & 9093 & Inhibition d'alertes via l'API publique & Retrait du binding hôte \\
        \hline
        VictoriaMetrics & 8428 & Lecture des métriques conservées sur 90 jours & Retrait du binding hôte (port déjà détecté par les scanners externes) \\
        \hline
        \texttt{devsecops-agent} & 8000 & Contournement de l'authentification Basic Auth de nginx & Retrait du binding hôte, accès exclusif via nginx \\
        \hline
        Open WebUI & 3001 & Création de compte ouverte (\texttt{ENABLE\_\allowbreak{}SIGNUP=true}) & Désactivation du \textit{signup} une fois le compte administrateur créé \\
        \hline
        \end{tabular}
        \caption{Écarts de durcissement identifiés sur la configuration de production}
        \label{tab:hardening_gaps}
        \end{table}

        La cause profonde est commune à toutes les lignes du tableau : ces \textit{bindings}
        ont été ajoutés un par un lors des phases de mise au point pour permettre
        l'inspection directe d'un service sans passer par le \textit{reverse proxy}, et
        n'ont pas été retirés une fois la phase terminée. Le correctif est symétrique :
        supprimer systématiquement les ports hôte, ne laisser que nginx exposé sur
        80 et 443, et accepter que tout diagnostic passe désormais soit par
        \texttt{docker exec} soit par une route nginx authentifiée. Le coût d'usage est
        marginal pour un développeur familier de la plateforme, et le gain de surface
        d'attaque est immédiat.

        Au-delà des ports, le durcissement complet suppose trois chantiers
        supplémentaires. Installer un certificat Let's Encrypt sur nginx chiffrerait
        tout le flux sortant du VPS et fermerait la dernière fenêtre en clair vers
        l'extérieur. Mettre en place une sauvegarde quotidienne de PostgreSQL, qui
        héberge à la fois la base de connaissances sécurité et les \textit{checkpoints}
        LangGraph, protégerait contre la perte du volume Docker ; la panne
        VictoriaMetrics d'avril a rappelé qu'un tel incident peut survenir sans prévenir.
        Restent enfin les sept conteneurs dépourvus de \textit{healthcheck} Docker
        (\texttt{nginx}, \texttt{prometheus}, \texttt{alertmanager}, \texttt{grafana},
        \texttt{victoriametrics}, \texttt{node-exporter} et \texttt{open-webui}) : les
        instrumenter laisserait le
        démon Docker détecter et relancer ces services de lui-même, sans attendre que
        les règles Prometheus finissent par déclencher une alerte. Ces trois chantiers
        prolongent la démarche engagée lors de l'audit du 28 avril et visent à rendre la
        plateforme robuste face aux incidents d'exploitation.

        Un dernier point, de nature différente, relève de la sécurité du moteur lui-même
        plutôt que de sa surface réseau. Le diff soumis au modèle 14B est un contenu
        rédigé par l'auteur de la \textit{Pull Request} analysée ; il constitue à ce titre
        une entrée non fiable. Un développeur malveillant pourrait y glisser des
        instructions destinées à détourner la revue, par exemple un commentaire de code
        enjoignant au modèle d'émettre un verdict \texttt{APPROVE}. Deux propriétés de
        l'architecture limitent la portée d'une telle tentative sans l'écarter
        complètement. D'une part, les cinq scanners sont déterministes et s'exécutent
        indépendamment du modèle : leurs \textit{findings} sont enregistrés dans
        PostgreSQL quel que soit le texte produit par le LLM, et un secret ou une CVE
        relevé par un scanner ne disparaît pas si le modèle est manipulé. D'autre part, la
        température fixée à 0,1 et l'extraction structurée du bloc JSON réduisent la
        latitude laissée au modèle dans la formulation de son verdict. La parade complète
        reste un chantier ouvert : isoler explicitement le diff par des délimiteurs dans le
        \textit{prompt}, et surtout faire reposer le \textit{gate} de blocage sur la
        sévérité des \textit{findings} déterministes des scanners plutôt que sur le seul
        verdict du modèle, de sorte qu'aucune injection dans le code revu ne puisse
        relâcher le contrôle de sécurité.

\section*{Conclusion}
\addcontentsline{toc}{section}{Conclusion}

Ce chapitre a retracé la réalisation effective du BTE Security AI Agent, conduite selon
la méthodologie Scrumban présentée au chapitre~1. Après avoir décrit l'environnement
matériel et logiciel d'accueil, nous avons suivi l'implémentation sprint par sprint, du
Sprint~1 (fondations Docker et webhook GitHub) au Sprint~8.1 (extension de
l'observabilité), en exposant à chaque itération l'objectif inscrit au backlog, les
obstacles rencontrés en production et le livrable obtenu. Certains épisodes ont davantage
marqué le projet : l'urgence disque d'avril, la remise en service des cibles Prometheus
restées \og down \fg{} après le premier déploiement, le travail sur node-exporter et la
règle iptables associée, la mise en place du garde-fou anti-hallucination du chat, la
réparation de l'AlertManager silencieusement hors service depuis le déploiement initial,
et la réécriture du parser de revue PR au Sprint~8. Aucun de ces problèmes ne figurait au
plan de départ : tous ont surgi en production avant d'être convertis en cartes prioritaires
du sprint en cours. C'est là le \textit{pull system} Scrumban à l'œuvre, chaque obstacle
remontant au backlog dès qu'il est constaté plutôt qu'au cycle de planification suivant. La
dernière section a rendu compte de la validation sur la PR de démonstration~\#18 : une revue
complète en une quinzaine de minutes, assortie de commentaires \textit{inline} et de
suggestions de correction publiés directement sur GitHub, la fourchette de quinze à
vingt-cinq minutes restant la cible pour des PR plus volumineuses.
